\chapter*{Abstract}
\addcontentsline{toc}{chapter}{Abstract}
\markboth{Abstract}{}

\noindent\textbf{Context.} This report is part of a two-month internship carried out
within the IT department of the Mövenpick Hôtel Tanger. Over time, a large number of
devices connect to the hotel's local network (workstations, printers, mobile terminals,
connected devices), without the IT department having a simple way to keep an up-to-date
inventory of them.

\vspace{0.3cm}
\noindent\textbf{Objective.} The objective of the project is to develop a network
device discovery and mapping tool, named \textbf{NetSentry}, able to scan the hotel's
local network, identify active devices (IP address, MAC address, vendor, open ports),
keep a history of scans, and automatically flag any unknown device appearing on the
network.

\vspace{0.3cm}
\noindent\textbf{Method.} The project relies on an initial upskilling phase covering
network fundamentals (IP/MAC addressing, the ARP protocol, ports), followed by the
design and development of a three-layer architecture: an Nmap-based scanning engine, a
backend API (Node.js/Express) handling orchestration, persistence (PostgreSQL) and
anomaly detection, and a web dashboard (React) allowing an IT administrator to view
devices, scan history and alerts. The solution is containerized with Docker to ease
deployment.

\vspace{0.3cm}
\noindent\textbf{Results.} The resulting tool scans the local network, lists connected
devices along with their vendor and open ports, keeps a full history of scans, and
raises an alert whenever an unrecognized device joins the network. The full chain --
scanning, backend, dashboard and containerized deployment -- was verified end to end
during the internship.

\vspace{0.5cm}
\noindent \textbf{Keywords:} Network discovery, Device mapping, Nmap, Anomaly
detection, Dashboard, Docker, NetSentry.
